Modern datacenter applications operate under strict microsecond-scale Service Level Objectives (SLOs) while exhibiting distinct workload characteristics. 
These workloads differ in request arrival patterns and service time distributions.
For example, microservice often experiences high request dispersion while serverless platforms and financial trading systems often exhibit more uniform workloads yet bursty traffic patterns.

The diversity in application workloads motivates the need for tailoring scheduling policies to meet specific SLOs.
Tail latency-sensitive services, such as in-memory key-value stores and web search, benefit from preemptive scheduling to minimize queueing delays. 
In contrast, interactive applications like e-commerce and social media prioritize responsiveness, often requiring policies like Shortest Remaining Processing Time (SRPT) to minimize mean latency.
Moreover, soft real-time workloads, including video conferencing and online gaming, operate under strict temporal constraints, demanding deadline-aware scheduling to guarantee frame rates and minimize jitter.

However, enforcing application-specific scheduling policies at microsecond-scale remains to be challenging.
Traditional kernel schedulers provide abundant scheduling policies but operate at millisecond timescales, unsuitable for modern applications.
To compensate, recent dataplane operating systems (OSes) bypass the kernel to achieve microsecond-scale preemptive scheduling in userspace.
Yet they lack the flexibility for applications to define custom scheduling policies.
Instead, they rigidly hardcode a hybrid policy by combining centralized First-Come-First-Served (FCFS) for light-tailed workloads with Processor Sharing (PS) for heavy-tailed workloads.

In this paper, we present the design of \kernel, a datapath OS that extends the Exokernel philosophy to hardware interrupt control.
Our key insight is that direct userspace access to system-critical interrupt hardware is the missing piece to achieving kernel-level scheduling flexibility in current dataplane OSes.
Therefore, the goal of \kernel is to redefine the OS-application boundary to safely expose these interrupts into userspace, thereby enabling applications to build custom scheduling policies directly around these low-level events.

\kernel relies on three key techniques to unite the policy flexibility of kernel schedulers with the microsecond-scale performance of dataplane OSes.
First, \kernel provides a scheduling core abstraction that enables applications to implement custom scheduling policies with kernel events in userspace.
Second, \kernel introduces a lightweight isolation mechanism that securely exposes low-level interrupt hardware to untrusted userspace components with minimal overhead.
Third, \kernel introduces a hardware multiplexing mechanism that safely shares system-critical resources (e.g., LAPIC timer) among multiple untrusted userspace schedulers.